\chapter{Potencia radiativa del fuego}\label{cap:frp}

\objetivos{
\begin{itemize}
  \item Definir la potencia radiativa del fuego (FRP) y explicar qué mide físicamente.
  \item Entender por qué puede estimarse con un solo canal de 4~\um.
  \item Relacionar la FRP con la biomasa consumida y con las emisiones.
  \item Explicar cómo se calcula en MODIS y en VIIRS y qué precauciones exige al comparar.
\end{itemize}}

\section{Más allá de contar}

Un foco dice <<aquí hay fuego>>. No dice cuánto. Una quema de rastrojo y el frente de un incendio
forestal generan cada uno un punto en el mapa. Para distinguirlos, los productos de fuego
incluyen una medida de intensidad.

\begin{definicion}{potencia radiativa del fuego (FRP)}
La potencia radiativa del fuego, o FRP por \emph{Fire Radiative Power}, es la energía que el
fuego contenido en un píxel emite en forma de radiación por unidad de tiempo, sumada sobre todas
las longitudes de onda y todas las direcciones. Se expresa en megavatios (MW).
\end{definicion}

Es la columna \texttt{frp} de los datos. Un megavatio es un millón de vatios; una quema agrícola
pequeña tiene una FRP de unos pocos MW, y los píxeles más intensos de un gran incendio forestal
superan los 1000~MW.

\section{De la ley de Stefan--Boltzmann a la FRP}

Si conociéramos el área $A_f$ y la temperatura $T_f$ del fuego, su potencia radiada saldría de
la ley de Stefan--Boltzmann, restando lo que esa misma superficie emitiría si estuviera a la
temperatura del fondo:
\begin{equation}\label{eq:frp-sb}
  \mathrm{FRP} = \sigma\,A_f\,\bigl(T_f^4 - T_b^4\bigr) \approx \sigma\,A_f\,T_f^4 .
\end{equation}
El término del fondo es despreciable: $300^4$ es menos del 1\,{}\% de $1000^4$.

\begin{ejemplo}{potencia de una quema pequeña}
Calcula la FRP de un fuego de \SI{100}{\metre\squared} a 1000~K y la de uno de
\SI{2000}{\metre\squared} a 800~K.
\solucion Para el primero: $\sigma A_f T_f^4 = \num{5.67e-8}\times 100\times 10^{12} =
\SI{5.67e6}{\watt} = \SI{5.7}{\mega\watt}$. Para el segundo: $\num{5.67e-8}\times 2000\times
(800)^4 = \num{5.67e-8}\times2000\times\num{4.096e11} = \SI{4.6e7}{\watt} =
\SI{46}{\mega\watt}$.
\end{ejemplo}

El problema es que ni $A_f$ ni $T_f$ se conocen: son justamente las incógnitas difíciles de
despejar del capítulo anterior. Se necesita un atajo.

\section{El atajo: en 4 micrómetros, la radiancia va como la cuarta potencia}

El atajo es una coincidencia afortunada de la ley de Planck. En el intervalo de temperaturas de
los fuegos de vegetación, entre unos 600 y 1500~K, la radiancia a 4~\um{} es casi exactamente
proporcional a $T^4$:
\begin{equation}\label{eq:aprox-mir}
  B_4(T) \approx a\,T^4 ,
\end{equation}
con una constante $a$ que depende de la banda concreta del sensor (para MODIS, $a \approx
\SI{3.0e-9}{\watt\per\square\metre\per\steradian\per\micro\metre\per\kelvin\tothe{4}}$).

La figura~\ref{fig:mir} lo comprueba. No es una ley exacta sino una aproximación válida en ese
intervalo, y funciona porque, para esas temperaturas, el canal de 4~\um{} está cerca del máximo
de la curva de Planck.

\begin{figure}[ht]
\centering
\begin{tikzpicture}
\begin{axis}[width=12.5cm,height=6.8cm,xmin=400,xmax=1600,ymin=0,ymax=2,
  xlabel={Temperatura del fuego (K)},
  ylabel={$B_4(T)\,/\,(a\,T^4)$},
  ytick={0,0.5,1,1.5,2},yticklabels={0,{0,5},1,{1,5},2},
  grid=major,grid style={ceniza!15},samples=120,domain=400:1600,
  declare function={planck(\l,\T)=1.191e8/(\l^5*(exp(14388/(\l*\T))-1));}]
\fill[bosque!15] (axis cs:600,0) rectangle (axis cs:1500,2);
\addplot[fuego,very thick] {planck(3.96,x)/(3.0e-9*x^4)};
\addplot[ceniza,dashed,domain=400:1600] {1};
\node[font=\footnotesize,bosque!70!black] at (axis cs:1050,1.75) {intervalo de los fuegos de vegetación};
\end{axis}
\end{tikzpicture}
\caption{Cociente entre la radiancia de Planck a 3,96~\um{} y la aproximación $a\,T^4$. Entre
600 y 1500~K (banda sombreada) el cociente se mantiene entre 0,74 y 1,07: con un error de
hasta un 25\,{}\% en los extremos, la radiancia del canal es proporcional a la potencia total
emitida.}
\label{fig:mir}
\end{figure}

Ahora combinemos tres ecuaciones. El exceso de radiancia del píxel sobre su fondo, según el
modelo del píxel mixto, es $\Delta L_4 = L_4 - L_{4,b} \approx p\,B_4(T_f)$. Con la
aproximación~\eqref{eq:aprox-mir}, $\Delta L_4 \approx p\,a\,T_f^4$. Y como $p =
A_f/A_{\text{píxel}}$:
\[
  \Delta L_4 \approx \frac{A_f}{A_{\text{píxel}}}\,a\,T_f^4
  \quad\Longrightarrow\quad
  A_f\,T_f^4 \approx \frac{A_{\text{píxel}}}{a}\,\Delta L_4 .
\]
Sustituyendo en la ecuación~\eqref{eq:frp-sb}:
\begin{equation}\label{eq:frp-mir}
  \boxed{\;\mathrm{FRP} \approx \frac{\sigma\,A_{\text{píxel}}}{a}\,\bigl(L_4 - L_{4,b}\bigr)\;}
\end{equation}

El producto desconocido $A_f T_f^4$ ha desaparecido. La potencia del fuego se obtiene de tres
cosas que sí se conocen: el área del píxel, la radiancia medida a 4~\um{} y la radiancia del
fondo estimada con los vecinos. Este es el \termino{método de la radiancia en el infrarrojo
medio}, debido a Wooster y colaboradores, y es el que usan los productos actuales.

\begin{ejemplo}{comprobación del método}
Verifica la ecuación~\eqref{eq:frp-mir} con el fuego de \SI{100}{\metre\squared} a 1000~K
dentro de un píxel de \SI{1}{\kilo\metre\squared}, usando los valores del capítulo anterior.
\solucion Allí se obtuvo $L_4 = 1{,}005$ con fondo $L_{4,b} = 0{,}673$, es decir, $\Delta L_4 =
0{,}332$ \si{\watt\per\square\metre\per\steradian\per\micro\metre}. Entonces
\[
 \mathrm{FRP} \approx \frac{\num{5.67e-8}\times 10^6}{\num{3.0e-9}}\times 0{,}332
 = \num{1.89e7}\times 0{,}332 = \SI{6.3e6}{\watt} = \SI{6.3}{\mega\watt}.
\]
El valor exacto del ejemplo anterior era 5,7~MW. La diferencia, de un 10\,{}\%, es el error de la
aproximación~\eqref{eq:aprox-mir} a 1000~K, visible en la figura~\ref{fig:mir}.
\end{ejemplo}

\section{Qué mide y qué no mide la FRP}

La FRP es una potencia: energía por unidad de tiempo en el instante de la pasada. Tres
precisiones.

\textbf{Es solo la parte radiada.} Un fuego libera energía también por convección (el aire
caliente que sube) y por conducción hacia el suelo. La fracción radiada es del orden de una
décima a una quinta parte del total.

\textbf{Es proporcional al ritmo de combustión.} Experimentos de quema controlada muestran que
la energía radiada total es proporcional a la masa de vegetación consumida, con un factor
cercano a 0,37~kg de biomasa seca por megajulio radiado. Por eso la FRP, integrada en el tiempo,
se usa para estimar cuánta biomasa se quemó y cuánto humo y CO$_2$ se emitió.

\textbf{Es instantánea.} Para pasar de potencia a energía hay que integrar sobre la duración del
fuego, y un satélite polar solo aporta unas pocas muestras al día. La integral, llamada
\termino{energía radiativa del fuego} (FRE), requiere suponer la forma del ciclo diurno.

\begin{ejemplo}{de potencia a biomasa}
Un píxel muestra una FRP de 40~MW. Suponiendo que el fuego mantiene esa potencia durante dos
horas, ¿cuánta biomasa consume?
\solucion La energía radiada es $E = \SI{40}{\mega\watt}\times\SI{7200}{\second} =
\SI{2.88e5}{\mega\joule}$. Con 0,37~kg por megajulio: $m = \num{2.88e5}\times 0{,}37 \approx
\SI{1.1e5}{\kilo\gram}$, unas 107 toneladas de materia seca. El resultado depende por completo
de la suposición sobre la duración.
\end{ejemplo}

\section{La FRP en MODIS y en VIIRS}

Los dos productos calculan la FRP con la ecuación~\eqref{eq:frp-mir}, pero con diferencias
importantes en los detalles.

\textbf{MODIS} usa su canal de 3,96~\um{} (banda 22, o 21 si la 22 se satura) y el área real del
píxel, que crece hacia los bordes. La versión actual del producto adoptó este método; las
versiones antiguas usaban una fórmula empírica distinta, basada en la octava potencia de la
temperatura de brillo.

\textbf{VIIRS de 375~m} no puede usar I4 porque se satura a 367~K. La detección se hace a
375~m, pero la FRP se calcula con la radiancia de la banda \textbf{M13}, de 750~m y con un rango
dinámico muy amplio, que no se satura con los fuegos de vegetación. La potencia obtenida para el
píxel de 750~m se reparte entre las detecciones de 375~m que contiene.

\begin{figure}[ht]
\centering
\begin{tikzpicture}[x=1.15cm,y=1.15cm]
  \fill[cenizaClaro] (0,0) rectangle (2,2);
  \draw[ceniza,very thick] (0,0) rectangle (2,2);
  \draw[ceniza] (1,0)--(1,2) (0,1)--(2,1);
  \fill[fuego!60] (0,1) rectangle (1,2); \fill[fuego!60] (1,1) rectangle (2,2); \fill[fuego!60] (1,0) rectangle (2,1);
  \draw[ceniza,very thick] (0,0) rectangle (2,2);
  \draw[ceniza] (1,0)--(1,2) (0,1)--(2,1);
  \node[font=\footnotesize] at (0.5,1.5) {foco}; \node[font=\footnotesize] at (1.5,1.5) {foco}; \node[font=\footnotesize] at (1.5,0.5) {foco};
  \draw[<->] (0,-0.25)--(2,-0.25) node[midway,below,font=\footnotesize]{píxel M13 de 750 m};
  \draw[<->] (2.25,1)--(2.25,2) node[midway,right,font=\footnotesize]{375 m};
  \node[caja,anchor=west,text width=7cm,align=left] at (4.2,1) {%
    FRP medida en el píxel de 750 m: 36 MW\\[3pt]
    Detecciones de 375 m dentro: 3\\[3pt]
    FRP asignada a cada foco: $36/3 = 12$ MW};
\end{tikzpicture}
\caption{En el producto de 375~m de VIIRS, la FRP se mide en la banda M13 de 750~m y se reparte
por igual entre los focos de 375~m contenidos en ese píxel.}
\label{fig:m13}
\end{figure}

\section{Precauciones al usar la FRP}

\begin{atencion}[La FRP de un foco no es comparable sin más entre sensores]
Hay al menos cuatro motivos.
\begin{enumerate}
  \item \textbf{Fuegos no detectados.} La FRP solo existe para los píxeles clasificados como
  fuego. MODIS deja fuera los fuegos pequeños y con ellos su potencia.
  \item \textbf{Reparto entre píxeles.} La potencia de un mismo frente se reparte entre más
  focos en VIIRS; cada foco tiene menos FRP pero hay más focos.
  \item \textbf{Geometría.} Hacia los bordes de la franja de MODIS, además de detectarse menos,
  la estimación se degrada.
  \item \textbf{Fondo.} La radiancia del fondo se estima con vecinos de tamaños muy distintos
  en cada sensor.
\end{enumerate}
Lo comparable no es la FRP de cada foco, sino la \emph{suma} de la FRP en una celda y un
periodo, y aun así con cautela.
\end{atencion}

\begin{parareto}
La FRP ofrece dos posibilidades para el proyecto. Como \textbf{segunda variable del calendario}:
además de <<cuántas detecciones>>, mostrar <<cuánta potencia>> permite distinguir una temporada
de muchas quemas pequeñas de una con pocos incendios muy intensos. Y como \textbf{criterio de
armonización}: descartar de VIIRS las detecciones con FRP por debajo de lo que MODIS podría
haber detectado es una forma físicamente motivada de acercar los dos registros
(capítulo~\ref{cap:armonizacion}). Empieza por los conteos; la FRP es una extensión.
\end{parareto}

\resumen

\begin{itemize}
  \item La \textbf{FRP} es la potencia que el fuego de un píxel emite como radiación, en MW.
  \item Por Stefan--Boltzmann, $\mathrm{FRP} \approx \sigma A_f T_f^4$, pero $A_f$ y $T_f$ no se
  conocen.
  \item Entre 600 y 1500~K la radiancia a 4~\um{} es proporcional a $T^4$, lo que permite
  calcular $\mathrm{FRP} \approx (\sigma A_{\text{píxel}}/a)\,(L_4 - L_{4,b})$ sin conocer el
  tamaño ni la temperatura del fuego.
  \item La FRP es proporcional al ritmo de consumo de biomasa; integrada en el tiempo da la
  energía radiada y, con ella, la masa quemada.
  \item MODIS la calcula con su canal de 3,96~\um; VIIRS de 375~m, con la banda M13 de 750~m,
  repartiéndola entre los focos.
  \item La FRP de focos individuales no es comparable entre sensores; las sumas por celda y
  periodo lo son algo más.
\end{itemize}

\ejercicios

\begin{enumerate}
  \item Calcula con la ecuación~\eqref{eq:frp-sb} la FRP de un frente de fuego de 300~m de largo
  y 5~m de profundidad a 900~K.
  \item Usando la figura~\ref{fig:mir}, ¿la ecuación~\eqref{eq:frp-mir} sobrestima o subestima
  la potencia de una brasa a 600~K? ¿Y la de una llama a 1200~K?
  \item Un píxel de MODIS con \texttt{scan} = 2,0 y \texttt{track} = 1,4 tiene un exceso de
  radiancia a 4~\um{} de 0,50 unidades sobre su fondo. Calcula su FRP.
  \item Un píxel M13 con 54~MW contiene cuatro detecciones de 375~m. ¿Qué FRP lleva cada una?
  Si MODIS observa el mismo fuego en un único píxel, ¿qué FRP esperarías, idealmente?
  \item Una región registra en un mes 400 focos con FRP media de 8~MW; otra, 60 focos con FRP
  media de 90~MW. Compara su FRP total y describe qué tipo de fuego sugiere cada una.
  \item Explica por qué la FRP medida a las 13:30 no puede multiplicarse por 24 horas para
  obtener la energía diaria.
\end{enumerate}

\lecturas

Wooster, Zhukov y Oertel (2003) derivan el método de la radiancia en el infrarrojo medio, y
Wooster y colaboradores (2005) establecen la relación entre energía radiada y biomasa
consumida. Las referencias están en el apéndice~\ref{ap:fuentes}.
